We presented a family of continuous hidden Markov models (CHMM-N, CHMM-t, CHMM-L) trained by EM for generating synthetic financial time series, together with two cross-asset dependence extensions for multi-asset synthesis.
All three variants share the same log-space forward-backward scaffold and quantile-based initialization; the M-step is the classical Baum-Welch update for Gaussian emissions, an ECM sequence with a golden-section search over $\nu_k$ for Student-t emissions, and closed-form weighted-median / weighted-MAD estimators for Laplace emissions.
The central univariate finding is that each CHMM variant reproduces all three canonical stylized facts of equity returns (heavy-tailed distributions, negligible linear autocorrelation, and persistent volatility clustering) at a single moderate state resolution ($K = 18$, selected from a $\{3, 6, 9, 12, 15, 18, 21\}$ sweep by winning or tying on every distributional pass-rate metric) \emph{without} requiring jump augmentation, discretization, or any mechanism beyond the standard HMM framework.
The heavy-tailed CHMM-t and CHMM-L variants additionally close the residual kurtosis gap of the Gaussian CHMM-N.

This result directly challenges the influential conclusion of \citet{ryden1998stylized} that Gaussian HMMs cannot reproduce the slow decay of the autocorrelation function of absolute returns.
We demonstrated that this limitation is an artifact of insufficient state resolution: \citet{ryden1998stylized} tested only $K = 2$--$3$ states, where the geometric sojourn-time distribution cannot approximate slow ACF decay.
At $K \geq 3$ with quantile-based initialization, the Baum-Welch algorithm learns a transition matrix with sufficient diagonal dominance in high-volatility states to sustain persistence, achieving ACF-MAE of 0.046--0.051 (comparable to GARCH(1,1) at 0.047) through a sum-of-exponentials approximation that effectively mimics the flexible sojourn-time distributions of hidden semi-Markov models \cite{bulla2006stylized}.

The direct head-to-head against the discrete framework of \citet{alswaidan2026hybrid}, reproduced as the Discrete NJ and Discrete WJ rows of Table~\ref{tab:model_comparison}, reveals a 94.7-percentage-point swing in IS KS pass rate (0\% for either discrete variant $\to$ 94.7\% for the CHMM-N).
The swing is attributable to the move from bin-centroid emissions to continuous Gaussian emissions: once the simulated return series carries continuous-valued information, standard distributional tests on the return series pass; once it is quantized to 13 centroids, they fail by construction.
The elimination of the jump mechanism further removes two hyperparameters ($\epsilon$ for jump probability, $\lambda$ for mean jump duration) and the associated grid search, reducing computational cost and model complexity.

The seven-model benchmark confirmed that no single generator dominates all metrics, but the CHMM at $K = 18$ achieves the most balanced performance: $94.7$--$95.2\%$ IS KS across the three emission families (vs.\ $23.1\%$ GARCH and $0.0\%$ for either discrete variant), ACF-MAE of $0.0513$ under CHMM-N (vs.\ $0.0492$ GARCH), Wasserstein-1 of $0.102$ under CHMM-t and CHMM-L (the smallest of any parametric model), and $100\%$ in-sample quantile coverage.
Cross-asset univariate generalization to NVDA ($96.7\%$ IS KS), JNJ ($99.4\%$), JPM ($97.8\%$), AAPL ($98.3\%$), and QQQ ($95.5\%$) confirmed adaptability to diverse equity risk profiles; out-of-sample evaluation on held-out 2024-01-04 through 2026-04-20 data showed robust generalization on SPY/JNJ/AAPL/QQQ but substantial OoS KS drops on NVDA ($55.8\%$) and JPM ($49.8\%$), which we attribute to the 2024--2026 macroeconomic regime (the AI-driven high-volatility regime for tech, and the interest-rate-repricing cycle for financials) and which the walk-forward refit partially recovers.

The cross-asset extensions then showed that the CHMM marginals compose cleanly with standard cross-sectional dependence models.
On the six-asset universe, both the Gaussian and Student-t copulas substantially outperformed the Single Index Model factor baseline: copulas achieved a median in-sample KS pass rate of $97.8\%$/$96.0\%$ versus $78.5\%$ for SIM, and an off-diagonal correlation MAE of $0.031$ (Gaussian) and $0.027$ (Student-t) versus $0.076$ for SIM.
The Student-t copula's profile MLE selected $\nu^* = 6$ degrees of freedom, and the resulting tail-dependent copula produced the best correlation reproduction overall.
These findings replicate the ordering reported by \citet{alswaidan2026hybrid} with discrete HMM marginals, suggesting that the copula advantage is intrinsic to the rank-reordering construction and transfers across marginal families.

Option pricing, implied-volatility-surface calibration, volatility-index generalization, and leverage-effect modelling are explicitly out of scope for this paper and are deferred to future work.

Several directions for future work follow from these results:
\begin{enumerate}
    \item \textit{Skew-heavy-tailed emissions.}
    This paper closes the kurtosis gap with Student-t and Laplace emissions; skew-$t$ or skew-Laplace emissions would additionally accommodate within-regime asymmetry and match the observed negative skewness of equity returns ($-0.75$ for SPY).

    \item \textit{Time-varying transition dynamics.}
    Estimating the transition matrix via rolling windows or Bayesian online learning would relax the stationarity assumption and should directly address the OoS cliff observed on NVDA and JPM over the $2024$--$2026$ OoS window.

    \item \textit{Full-universe multi-asset generation via vine copulas.}
    Scaling the copula construction from six assets to the 424-asset universe of \citet{alswaidan2026hybrid} requires vine or factor copulas for tractability and is a natural next step.

    \item \textit{Principled state selection.}
    Our $K = 18$ choice is a multi-metric score on top of a finite sweep; combining it with BIC or held-out log-likelihood (walk-forward) would yield a formally principled selection that plugs directly into the appendix Table~\ref{tab:sensitivity}.

    \item \textit{Regime-aware dynamic asset allocation.}
    The most immediate downstream application of the synthetic-data engine is regime-switching portfolio choice in the spirit of \citet{baekimmulvey2014}: using CHMM-decoded states as conditioning variables in a dynamic allocation model, and evaluating the realized utility of the resulting policy against policies calibrated on i.i.d.\ or GARCH generators.
\end{enumerate}

\paragraph{Data and Code Availability.}
The simulation scripts, training and test data, and all code to reproduce the results in this paper are available under an MIT license from the paper repository: \url{https://github.com/altashly1/CHMM-paper}.
The core continuous HMM implementation (all three emission families, SIM, and Gaussian / Student-t copula modules) is available in the code repository: \url{https://github.com/altashly1/CHMM-Model}.
All experiments reported in Section~\ref{sec:results} and in the appendix use a deterministic global seed ($20260420$), with per-experiment sub-seeds derived additively to keep each block independently reproducible.
The software environment is Julia $\geq 1.12$ with package versions pinned in \texttt{Manifest.toml}; running \texttt{Pkg.instantiate()} inside the project directory reproduces the dependency graph, after which the analysis entry-point script regenerates every table and figure.

\paragraph{Conflict of Interest.}
The authors declare no competing interests.

\paragraph{Author Contributions.}
J.V.\ directed the study.
A.A.\ developed the continuous model, implemented the Baum-Welch algorithm, implemented the Single Index Model and copula-based cross-asset extensions, conducted the empirical analysis, and generated all figures and tables.
C.J.\ contributed to methodology review and validation.
All authors edited and reviewed the final manuscript.
